\section{\label{sec:influence_to_hmf}From Influence Functionals to the Hamiltonian of Mean Force}
Having established the equivalence between the quenched density and the standard Euclidean path integral in the previous section, we now outline the specific form of the reduced system when coupled to a harmonic bath. Following the influence-functional formalism \cite{feynmanTheoryGeneralQuantum1963a,caldeiraQuantumTunnellingDissipative1983a,grabertQuantumBrownianMotion1988} and specifically its stochastic representation for Gaussian baths \cite{mccaulPartitionfreeApproachOpen2017c, mccaulDrivingSpinbosonModels2018a, mccaulHowWinFriends2021b}, the resulting reduced density matrix admits the imaginary-time path integral representation:
\begin{widetext}
\begin{equation}
    \bar\rho_Q(q_f,q_i;\beta) = Z_X(\beta) \int_{q(0)=q_i}^{q(\beta)=q_f}\!\!\mathcal D q\; \exp\Bigg[ -S_Q[q] + \frac{1}{2}\int_0^\beta d\tau\int_0^\beta d\tau' f\big(q(\tau)\big) K(\tau-\tau') f\big(q(\tau')\big) \Bigg].
    \label{eq:reduced_effective_PI}
\end{equation}
\end{widetext}
Here $Z_X(\beta)=\Tr_X e^{-\beta H_X}$ is the free bath partition function, $S_Q$ is the Euclidean system action defined and $K(\tau)$ is the imaginary-time influence kernel. In terms of the bath spectral density $J(\omega)$, the kernel is given by~\cite{weissQuantumDissipativeSystems2012, feynmanTheoryGeneralQuantum1963a}
\begin{equation}
    K(\tau) = \frac{1}{\pi}\int_0^\infty d\omega\, J(\omega) \frac{\cosh\!\left[\omega\left(\frac{\beta}{2}-|\tau|\right)\right]} {\sinh\!\left(\frac{\beta\omega}{2}\right)}.
    \label{eq:kernel_spectral}
\end{equation}

While the representation \eqref{eq:reduced_effective_PI} is exact, it is nonlocal (bilocal) in imaginary time. This nonlocality can be resolved by a Hubbard - Stratonovich (HS) transformation \cite{hubbardCalculationPartitionFunctions1959a, stratonovich1957QDistro}, which decouples the bilocal interaction by introducing an auxiliary Gaussian field $\xi(\tau)$. This step provides the bridge to an operator-level representation of $\bar{\rho}_Q$ as a Gaussian average over the local quenched canonical densities discussed in Sec.~\ref{sec:quenched}:
\begin{equation}
\begin{split}
    \bar\rho_Q(\beta) 
    &= 
    Z_X(\beta)\,\mathbb E_\xi\bigg[ \\
    &\quad \hat{\mathcal{T}}\exp\!\left(-\int_0^\beta d\tau\,\bar{H}(\tau;\xi)\right) \bigg],
\end{split}
    \label{eq:bar_rho_as_quenched_average}
\end{equation}
The statistical properties of the auxiliary field $\xi(\tau)$ are set by the influence kernel \cite{mccaulPartitionfreeApproachOpen2017c,mccaulDrivingSpinbosonModels2018a,mccaulHowWinFriends2021b}
\begin{equation}
    \mathbb{E}_\xi[\xi(\tau)] = 0, \quad \mathbb{E}_\xi[\xi(\tau)\xi(\tau')] = K(\tau-\tau').
    \label{eq:noise_statistics}
\end{equation}
This operator representation maps directly to the quenched density framework. The challenge now lies in performing the average over $\xi$. To do so, we must isolate the stochastic term. Moving to the interaction picture with respect to $H_Q$, the propagator factorises as
\begin{equation}
    \hat{\mathcal{T}}e^{-\int_0^\beta d\tau\,[H_Q - \xi(\tau)f]} = e^{-\beta H_Q} \hat{\mathcal{T}}e^{\int_0^\beta d\tau\, \xi(\tau) f(\tau)},
    \label{eq:interaction_picture_split}
\end{equation}
where $f(\tau) = e^{\tau H_Q} f e^{-\tau H_Q}$. While we have successfully isolated the stochastic field $\xi(\tau)$, the coupling operator $f(\tau)$ has acquired a nontrivial imaginary-time dependence. This greatly complicates the averaging out of the stochastic term, and we defer its explicit evaluation to future work. 

Here we focus instead on the commuting (QND) sector $[H_Q,f]=0$, where the algebraic obstruction disappears: $f(\tau)=f$ is $\tau$-independent, the $\hat{\mathcal{T}}$ ordering becomes redundant and the expression factorises:
\begin{align}
    \hat{\mathcal{T}}\exp\!\left[-\int_0^\beta d\tau\,\bar{H}(\tau;\xi)\right]
    &=
    e^{-\beta H_Q}\, \exp\!\left[f\,X\right],
    \label{eq:tilde_rho_commuting_factorised}
\end{align}
where $X:=\int_0^\beta d\tau\,\xi(\tau)$ is a Gaussian variable with mean zero and variance
\begin{equation}
    C(\beta) = \mathbb E_\xi[X^2] = \int_0^\beta d\tau\int_0^\beta d\tau'\,K(\tau-\tau').
    \label{eq:Cbeta_def}
\end{equation}
Substituting the kernel Eq.~\eqref{eq:kernel_spectral} and performing the double
imaginary-time integral gives
\begin{align}
    C(\beta)
    &=\frac{2\beta}{\pi}\int_0^\infty d\omega\,\frac{J(\omega)}{\omega}
    \equiv 2\beta\,\lambda,
    \nonumber\\
    \lambda&:=\frac{1}{\pi}\int_0^\infty d\omega\,\frac{J(\omega)}{\omega},
    \label{eq:Cbeta_explicit}
\end{align}
where $\lambda$ is the reorganisation energy~\cite{weissQuantumDissipativeSystems2012, leggettDynamicsDissipativeTwostate1987}.
Crucially, $C(\beta)$ is \emph{linear} in $\beta$, so the ratio $C(\beta)/(2\beta)=\lambda$ is temperature independent.

With this the Gaussian average $\mathbb E_\xi[e^{f X}] = \exp[\frac{1}{2}C(\beta)f^2]$ can be performed exactly, and using Eq.~\eqref{eq:bar_rho_as_quenched_average}, we obtain:
\begin{equation}
    \bar\rho_Q(\beta) = Z_X(\beta)\,e^{-\beta H_Q}\,\exp\!\left[\beta\lambda\,f^2\right].
    \label{eq:bar_rho_commuting_closed}
\end{equation}

This is (up to normalisation) the canonical exponential form for a thermal state, and the Hamiltonian of mean force is identically
\begin{equation}
    H_{\mathrm{MF}}(\beta) = H_Q -\lambda\,f^2 
    \label{eq:HMF_commuting_final}
\end{equation}
up to some temperature dependent constant that is fixed by the normalisation
$\Tr\,e^{-\beta H_{\mathrm{MF}}}=1$. The only nontrivial operator content is then the correction $-\lambda f^2$, which is temperature independent.

Interestingly $\lambda$ coincides with the static potential renormalisation
$\frac{1}{2}\sum_k c_k^2/(m_k\omega_k^2)$ of the Caldeira-Leggett model~\cite{correaPotentialRenormalisationLamb2025}. If the standard counterterm $+\lambda f^2$ has been absorbed into $H_Q$ (as stated in Sec.~\ref{sec:model}), the bath-induced shift
$-\lambda f^2$ cancels it exactly. In this case, $H_{\mathrm{MF}}$ reduces to the bare system Hamiltonian $H_Q^{\mathrm{bare}}$, and the normalised mean-force Gibbs state coincides with the bare-system Gibbs state. In other words, in the commuting sector with translational invariance enforced, the bath has no operator-level effect on the reduced equilibrium
state. Conversely, if no counterterm is included, $-\lambda f^2$ represents a genuine
bath-induced static renormalisation of the system potential.

Having established the quenched density representation exactly captures the reorganisation-energy result~\cite{campisiTalknerHanggi2009Solvable}, we now apply it to stochastic numerical evaluation of the equilibrium state.

